\section{Methodology}
\label{sec:methodology}

\subsection{Notation}
\label{sec:notation}

Let $G = (V, E, w)$ be the census-tract adjacency graph for a state, where
$V$ is the set of census tracts, $E$ is the set of shared-boundary edges, and
$w: E \to \mathbb{R}_{>0}$ assigns each edge the TIGER boundary length between
adjacent tracts. Let $\text{pop}: V \to \mathbb{Z}_{>0}$ be the 2020 census
population of each tract. Let $n$ be the state's apportioned congressional seat
count.

\subsection{The Split Prescription}
\label{sec:factorization}

At each level of the AR tree, a region with $k$ target districts is split
according to the following rule with threshold $\kappa = 3$:

\begin{definition}[Split prescription]
  \label{def:split}
  For $k \geq 1$, $\text{split}(k)$ is:
  \begin{enumerate}
    \item \textbf{Direct} ($k \leq \kappa$): split into $k$ equal parts
          (sub-regions each with target~1).
    \item \textbf{Uniform-composite} ($k > \kappa$, $k$ composite): let $p$ be
          the \emph{largest} prime factor of $k$. Split into $p$ equal parts,
          each with target $k/p$.
    \item \textbf{Binary fallback} ($k > \kappa$, $k$ prime): split into 2 parts
          with population targets $\lfloor k/2 \rfloor$ and $\lceil k/2 \rceil$.
  \end{enumerate}
\end{definition}

Examples: $\text{split}(8) = (2, \text{sub}=4)$; $\text{split}(9) = (3, \text{sub}=3)$;
$\text{split}(17) = (2, \text{sub}=(8,9))$; $\text{split}(34) = (17, \text{sub}=2)$;
$\text{split}(51) = (17, \text{sub}=3)$.

\paragraph{Why largest prime first?} The reuse property requires that two seat
counts $n$ and $n'$ sharing a common factor $p$ produce the same top-level
$p$-way partition. With largest-prime-first, $n=34=2\times 17$ and
$n'=51=3\times 17$ both prescribe a 17-way primary split, sharing the same
17-district top-level partition of any given state. Smallest-prime-first would
instead bisect the state first (different halves for $n=34$ vs.\ thirds for
$n'=51$), destroying reuse.

\paragraph{Why binary fallback for large primes?} A direct $k$-way METIS
cut for large prime $k$ (e.g.\ $k=17$) is computationally feasible but
produces no reuse with neighbouring seat counts and is harder to balance than
iterated bisections. The binary fallback $17 \to (8, 9)$ naturally decomposes
into the highly-reused sub-trees $k=8=2^3$ (three binary levels) and $k=9=3^2$
(two ternary levels).

\subsection{The AR Algorithm}
\label{sec:pfr}

\begin{definition}[ApportionRegions]
  Given graph $G$ and seat count $n$, the AR map $M_n(G)$ is constructed by
  the recursive procedure $\text{AR}(G, n)$:
  \begin{enumerate}
    \item If $n = 1$: return $\{G\}$ (one district = whole region).
    \item Compute $\text{split}(n) = (p, k_1, \ldots, k_p)$ where $p$ is the
          number of parts and $k_i$ is the target district count for part $i$
          (Definition~\ref{def:split}).
    \item Partition $G$ into $p$ sub-regions $R_1, \ldots, R_p$ with target
          population fractions $k_i / n$ each, minimising edge-cut weight.
    \item For each $i$: recursively compute $\text{AR}(R_i, k_i)$.
    \item Return the union of all recursive results.
  \end{enumerate}
\end{definition}

The $p$-way partition in step~3 is computed by METIS~5.x (version with
\texttt{part\_kway} for $p > 2$, \texttt{part\_recursive} for $p = 2$) using
population vertex weights and TIGER boundary-length edge weights. For the
uniform case ($k_1 = \cdots = k_p = n/p$) METIS uses equal target fractions;
for the binary fallback case ($k_1 = \lfloor n/2 \rfloor$, $k_2 = \lceil n/2
\rceil$) it uses proportional target fractions $k_1/n$ and $k_2/n$.

\paragraph{Determinism assumption.}
METIS~5.x's multilevel $k$-way partitioning is deterministic given a fixed
random seed and fixed input graph (adjacency structure, vertex weights, edge
weights). We rely on this property throughout: when we say two seat counts
``share the same partition,'' we mean the assignment of census tracts to
districts is byte-identical, conditional on the same METIS seed at that level.
Our implementation sets the per-level seed as a deterministic function of the
global content-derived seed and the tree depth, ensuring full reproducibility.

\subsection{The Reuse Property}
\label{sec:reuse}

\begin{definition}[Canonical $p$-partition of $G$]
  Fix a census-tract graph $G$ and a positive integer $p \leq |V(G)|$. The
  \emph{canonical $p$-partition} $\Pi_p(G, \sigma)$ is the output of METIS
  $k$-way partitioning on $G$ with $k = p$, equal population target fractions,
  and METIS seed $\sigma$. By the determinism assumption, $\Pi_p(G, \sigma)$ is
  a function of $(G, p, \sigma)$ alone.
\end{definition}

\begin{theorem}[Reuse]
  \label{thm:reuse}
  Let $n$ and $n'$ be composite integers with the same largest prime factor $p$:
  $n = p \cdot m$, $n' = p \cdot m'$, with $p > \text{lpf}(m)$ and
  $p > \text{lpf}(m')$ (where $\text{lpf}$ denotes the largest prime factor).
  Let $\sigma$ be any METIS seed and $G$ any census-tract graph. Then
  $\text{AR}(G, n, \sigma)$ and $\text{AR}(G, n', \sigma)$ begin with the
  same top-level partition $\Pi_p(G, \sigma_0)$, where $\sigma_0$ is the
  depth-0 seed derived from $\sigma$.
\end{theorem}

\begin{proof}
  Since $p = \text{lpf}(n) = \text{lpf}(n')$, Definition~\ref{def:split}
  case~(2) applies to both. At depth~0, both algorithms call METIS with
  inputs $(G, p, \sigma_0)$ and equal target fractions $1/p$. By the
  determinism assumption, both calls return $\Pi_p(G, \sigma_0)$
  identically. The sub-regions $R_1, \ldots, R_p$ are therefore identical
  for both seat counts; the two algorithms diverge only in the sub-district
  counts $(k_i = m)$ vs.\ $(k_i' = m')$ passed to the recursive calls.
  \qed
\end{proof}

\begin{corollary}
  $\text{AR}(G, 34, \sigma)$ and $\text{AR}(G, 51, \sigma)$ begin with
  $\Pi_{17}(G, \sigma_0)$: the same 17-district partition of $G$.
  After the shared first level, $k=34$ bisects each region while
  $k=51$ trisects each region. The partition cache in \texttt{PfrCompositor}
  is a direct implementation of this theorem: a cache hit on
  $(\text{hash}(G), 17, \sigma_0)$ serves both computations.
\end{corollary}

\paragraph{Partition cache correctness.}
The cache key $(\text{hash}(G), p, \sigma_0)$ uses a content-derived hash
of the census-tract graph (adjacency list + population vector). This is
content-addressed, not identifier-addressed: two subgraphs with identical
adjacency and population produce the same cache key regardless of their
position in the factorization tree. Across decennial reapportionments,
the hash changes with the census data, so the cache correctly invalidates
stale entries.

\paragraph{Formal status of the binary fallback.}
For prime $k > \kappa$ (threshold $\kappa = 3$), the split prescription uses
a 2-way METIS call with asymmetric population targets $\lfloor k/2 \rfloor$
and $\lceil k/2 \rceil$ rather than a direct $k$-way call. We label this
\emph{AR-binary}. AR-binary is formally distinct from direct-$k$-METIS:
the two partitions are outputs of different METIS invocations and may produce
different assignments. However, both are computed from the same inputs and
seed, so the Reuse Theorem still applies: two seat counts $n, n'$ sharing the
same large prime $p$ both call 2-way METIS with targets $(\lfloor p/2 \rfloor,
\lceil p/2 \rceil)$ on the same whole-state graph, returning the same binary
split regardless of whether $p$ was intended as a ``$p$-way cut.'' The reuse
theorem holds for AR-binary under the same determinism assumption.

A variant \emph{AR-direct} that uses direct $k$-way METIS for all prime $k$
is fully equivalent to AR-binary in our implementation: METIS~5.x's
\texttt{part\_kway} handles $k = 17$ directly (single call, 0.1\% tolerance).
We defer a systematic comparison of AR-binary vs.\ AR-direct to future work.

\subsection{Factorization Tree Structure}
\label{sec:tree}

The factorization tree $T_n$ for a state has:
\begin{itemize}
  \item \textbf{Root}: the whole state.
  \item \textbf{Level $\ell$ nodes}: regions produced by applying $q_1, \ldots,
        q_\ell$ to the root. Each level-$\ell$ node has $q_\ell$ children.
  \item \textbf{Leaves}: the $n$ districts.
  \item \textbf{Depth}: $\Omega(n)$ (number of prime factors with multiplicity).
\end{itemize}

For $n = 2^k$ (e.g., Wisconsin: $n = 8 = 2^3$), the tree is a complete binary
tree of depth~$k$ --- exactly the recursive bisection algorithm of
\citet{dellaLibera2026mec}.

For $n$ prime (e.g., Pennsylvania: $n = 17$), the tree has depth~1: a single
$n$-way cut of the whole state. No reuse is possible across different prime seat
counts.

\subsection{Handling Reapportionment}
\label{sec:reapportionment}

When $n$ changes to $n' = n \pm 1$ after a decennial census, the required
recomputation depends on the first position where $F(n)$ and $F(n')$ differ.

\begin{proposition}[Recomputation depth]
  Let $F(n) = [q_1, \ldots, q_d]$ and $F(n') = [q_1', \ldots, q_{d'}']$.
  Let $k^* = \min\{k : q_k \neq q_k'\}$ (or $k^* = 1$ if $q_1 \neq q_1'$).
  Then AR reuses levels $1, \ldots, k^*-1$ and recomputes levels
  $k^*, \ldots, \max(d, d')$.
\end{proposition}

The worst case is $k^* = 1$: the largest prime factor changes, requiring a
full top-level recomputation. This occurs for the $51 \to 52$ transition:
$51 = 3 \times 17$ (largest prime 17) vs.\ $52 = 4 \times 13$ (largest prime 13).
No geographic sub-region is shared between the two maps.

\subsection{AR-Smooth: Minimising Reapportionment Disruption}
\label{sec:smooth}

When the prime factorization changes structurally, AR-smooth finds the
largest common prefix of $F(n)$ and $F(n')$ and reuses it, completing the
remaining levels from the best available cached splits. Specifically, for
$n \to n+1$:
\begin{enumerate}
  \item Find the longest common prefix of $F(n)$ and $F(n+1)$.
  \item Reuse those levels' cached splits.
  \item Re-solve only the subtrees that changed.
\end{enumerate}
This minimises the number of census tracts that change district assignment
across reapportionments.

\subsection{Reproducibility}
\label{sec:methodology:reproducibility}

All AR runs in this paper use METIS~5.2 (vendored via the \texttt{redist-metis}
crate v0.2, statically linked).
METIS parameters: \texttt{ufactor}=5 (0.5\,\% per-part population balance
tolerance); \texttt{niter}=100 refinement iterations; \texttt{ncuts}=1
(single cut attempt per invocation, determinism enforced by fixed seed).
The global seed for each state is derived as:
\[
  \sigma_{\text{global}} = \mathrm{SHA\text{-}256}(\texttt{census\_release\_id}
    \mathbin{\|} \texttt{"DIA\_SEED\_V1"} \mathbin{\|} 0) \bmod 2^{31}
\]
where \texttt{census\_release\_id} is the Census Bureau's formal identifier for
the 2020 Redistricting Data release (P.L.~94-171, release date 2021-08-12).
Per-level seeds are derived deterministically from $\sigma_{\text{global}}$ and
the tree depth (bitwise rotation), ensuring full reproducibility without manual
seed specification.
The complete source code and build manifest (Rust workspace lockfile, METIS
vendored source, and pipeline configuration) are available at the project
repository.
